\documentclass[11pt,a4paper,oneside]{report}             % Single-side
%\documentclass[11pt,a4paper,twoside,openright]{report}  % Duplex

%\PassOptionsToPackage{chapternumber=Huordinal}{magyar.ldf}
\usepackage{t1enc}
\usepackage[latin2]{inputenc}
\usepackage{amsmath}
\usepackage{amssymb}
\usepackage{enumerate}
\usepackage[thmmarks]{ntheorem}
\usepackage{graphics}
\usepackage{epsfig}
\usepackage{listings}
\usepackage{color}
%\usepackage{fancyhdr}
\usepackage{lastpage}
\usepackage{anysize}
\usepackage[magyar]{babel}
\usepackage{sectsty}
\usepackage{setspace}  % Ettol a tablazatok, abrak, labjegyzetek maradnak 1-es sorkozzel!
\usepackage[hang]{caption}
\usepackage{hyperref}

%--------------------------------------------------------------------------------------
% Main variables
%--------------------------------------------------------------------------------------
\newcommand{\vikszerzo}{B�dis-Szomor� Andr�s}
\newcommand{\vikkonzulens}{dr.~Konzulens Elem�r}
\newcommand{\vikcim}{Elektronikus terel�k}
\newcommand{\viktanszek}{M�r�stechnika �s Inform�ci�s Rendszerek Tansz�k}
\newcommand{\vikdoktipus}{Diplomaterv}
\newcommand{\vikdepartmentr}{B�dis-Szomor� Andr�s}

%--------------------------------------------------------------------------------------
% Page layout setup
%--------------------------------------------------------------------------------------
% we need to redefine the pagestyle plain
% another possibility is to use the body of this command without \fancypagestyle
% and use \pagestyle{fancy} but in that case the special pages
% (like the ToC, the References, and the Chapter pages)remain in plane style

\pagestyle{plain}
%\setlength{\parindent}{0pt} % �ttekinthet�bb, angol nyelv� dokumentumokban jellemz�
%\setlength{\parskip}{8pt plus 3pt minus 3pt} % �ttekinthet�bb, angol nyelv� dokumentumokban jellemz�
\setlength{\parindent}{12pt} % magyar nyelv� dokumentumokban jellemz�
\setlength{\parskip}{0pt}    % magyar nyelv� dokumentumokban jellemz�

\marginsize{35mm}{25mm}{15mm}{15mm} % anysize package
\setcounter{secnumdepth}{0}
\sectionfont{\large\upshape\bfseries}
\setcounter{secnumdepth}{2}
\singlespacing
\frenchspacing

%--------------------------------------------------------------------------------------
% Setup hyperref package
%--------------------------------------------------------------------------------------
\hypersetup{
    bookmarks=true,            % show bookmarks bar?
    unicode=false,             % non-Latin characters in Acrobat�s bookmarks
    pdftitle={\vikcim},        % title
    pdfauthor={\vikszerzo},    % author
    pdfsubject={\vikdoktipus}, % subject of the document
    pdfcreator={\vikszerzo},   % creator of the document
    pdfproducer={Producer},    % producer of the document
    pdfkeywords={keywords},    % list of keywords
    pdfnewwindow=true,         % links in new window
    colorlinks=true,           % false: boxed links; true: colored links
    linkcolor=black,           % color of internal links
    citecolor=black,           % color of links to bibliography
    filecolor=black,           % color of file links
    urlcolor=black             % color of external links
}

%--------------------------------------------------------------------------------------
% Set up listings
%--------------------------------------------------------------------------------------
\lstset{
    basicstyle=\scriptsize\ttfamily, % print whole listing small
    keywordstyle=\color{black}\bfseries\underbar, % underlined bold black keywords
    identifierstyle=,           % nothing happens
    commentstyle=\color{white}, % white comments
    stringstyle=\scriptsize\sffamily, % typewriter type for strings
    showstringspaces=false,     % no special string spaces
    aboveskip=3pt,
    belowskip=3pt,
    columns=fixed,
    backgroundcolor=\color{lightgray},
}
\def\lstlistingname{lista}

%--------------------------------------------------------------------------------------
% Some new commands and declarations
%--------------------------------------------------------------------------------------
\newcommand{\code}[1]{{\upshape\ttfamily\scriptsize\indent #1}}

% define references
\newcommand{\figref}[1]{\ref{fig:#1}.}
\renewcommand{\eqref}[1]{(\ref{eq:#1})}
\newcommand{\listref}[1]{\ref{listing:#1}.}
\newcommand{\sectref}[1]{\ref{sect:#1}}
\newcommand{\tabref}[1]{\ref{tab:#1}.}

\DeclareMathOperator*{\argmax}{arg\,max}
%\DeclareMathOperator*[1]{\floor}{arg\,max}
\DeclareMathOperator{\sign}{sgn}
\DeclareMathOperator{\rot}{rot}
\definecolor{lightgray}{rgb}{0.95,0.95,0.95}

\author{\vikszerzo}
\title{\viktitle}
\includeonly{
    project,%
    titlepage,%
    declaration,%
    abstract,%
    introduction,%
    chapter1,%
    acknowledgement,%
    appendices,%
}
%--------------------------------------------------------------------------------------
% Setup captions
%--------------------------------------------------------------------------------------
\captionsetup[figure]{
%labelsep=none,
%font={footnotesize,it},
%justification=justified,
width=.75\textwidth,
aboveskip=10pt}

\renewcommand{\captionlabelfont}{\small\bf}
\renewcommand{\captionfont}{\footnotesize\it}

%--------------------------------------------------------------------------------------
% Table of contents and the main text
%--------------------------------------------------------------------------------------
\begin{document}
\singlespacing
%--------------------------------------------------------------------------------------
% Feladatkiiras (a tanszeken atveheto, kinyomtatott valtozat)
%--------------------------------------------------------------------------------------
\clearpage
\begin{center}
\large
\textbf{FELADATKI�R�S}\\
\end{center}

A feladatki�r�st a tansz�ki adminisztr�ci�ban lehet �tvenni, �s a leadott munk�ba eredeti, tansz�ki pecs�ttel ell�tott �s a tansz�kvezet� �ltal al��rt lapot kell belef�zni (ezen oldal \emph{helyett}, ez az oldal csak �tmutat�s). Az elektronikusan felt�lt�tt dolgozatban m�r nem kell beleszerkeszteni ezt a feladatki�r�st.


\pagenumbering{arabic}
\onehalfspacing
%--------------------------------------------------------------------------------------
% The title page
%--------------------------------------------------------------------------------------
\begin{titlepage}
\begin{center}
\includegraphics[width=60mm,keepaspectratio]{figures/BMElogo.png}\\
\vspace{0.3cm}
\textbf{Budapesti M�szaki �s Gazdas�gtudom�nyi Egyetem}\\
\textmd{Villamosm�rn�ki �s Informatikai Kar}\\
\textmd{\viktanszek}\\[5cm]

\vspace{0.4cm}
{\huge \bfseries \vikcim}\\[0.8cm]
\vspace{0.5cm}
\textsc{\Large \vikdoktipus}\\[4cm]

\begin{tabular}{cc}
 \makebox[7cm]{\emph{K�sz�tette}} & \makebox[7cm]{\emph{Konzulens}} \\
 \makebox[7cm]{\vikszerzo} & \makebox[7cm]{\vikkonzulens}
\end{tabular}

\vfill
{\large \today}
\end{center}
\end{titlepage}

\tableofcontents\vfill
%--------------------------------------------------------------------------------------
% Nyilatkozat
%--------------------------------------------------------------------------------------
\begin{center}
\large
\textbf{HALLGAT�I NYILATKOZAT}\\
\end{center}

Alul�rott \emph{\vikszerzo}, szigorl� hallgat� kijelentem, hogy ezt a szakdolgozatot/ diplomatervet \textcolor{blue}{(nem k�v�nt t�rlend�)} meg nem engedett seg�ts�g n�lk�l, saj�t magam k�sz�tettem, csak a megadott forr�sokat (szakirodalom, eszk�z�k stb.) haszn�ltam fel. Minden olyan r�szt, melyet sz� szerint, vagy azonos �rtelemben, de �tfogalmazva m�s forr�sb�l �tvettem, egy�rtelm�en, a forr�s megad�s�val megjel�ltem.

Hozz�j�rulok, hogy a jelen munk�m alapadatait (szerz�(k), c�m, angol �s magyar nyelv� tartalmi kivonat, k�sz�t�s �ve, konzulens(ek) neve) a BME VIK nyilv�nosan hozz�f�rhet� elektronikus form�ban, a munka teljes sz�veg�t pedig az egyetem bels� h�l�zat�n kereszt�l (vagy autentik�lt felhaszn�l�k sz�m�ra) k�zz�tegye. Kijelentem, hogy a beny�jtott munka �s annak elektronikus verzi�ja megegyezik. D�k�ni enged�llyel titkos�tott diplomatervek eset�n a dolgozat sz�vege csak 3 �v eltelte ut�n v�lik hozz�f�rhet�v�.

\begin{flushleft}
\vspace*{1cm}
Budapest, \today
\end{flushleft}

\begin{flushright}
 \vspace*{1cm}
 \makebox[7cm]{\rule{6cm}{.4pt}}\\
 \makebox[7cm]{\emph{\vikszerzo}}\\
 \makebox[7cm]{hallgat�}
\end{flushright}
\thispagestyle{empty}

\vfill
\clearpage
\thispagestyle{empty} % an empty page

%----------------------------------------------------------------------------
% Abstract in hungarian
%----------------------------------------------------------------------------
\chapter*{Kivonat}\addcontentsline{toc}{chapter}{Kivonat}

Jelen dokumentum egy diplomaterv sablon, amely formai keretet ad a BME Villamosm�rn�ki �s Informatikai Kar�n v�gz� hallgat�k �ltal elk�sz�tend� szakdolgozatnak �s diplomatervnek. A sablon haszn�lata opcion�lis. Ez a sablon \LaTeX~alap�, a \emph{TeXLive} \TeX-implement�ci�val �s a PDF-\LaTeX~ford�t�val m�k�d�k�pes.
\vfill

%----------------------------------------------------------------------------
% Abstract in english
%----------------------------------------------------------------------------
\chapter*{Abstract}\addcontentsline{toc}{chapter}{Abstract}

This document is a \LaTeX-based skeleton for BSc/MSc~theses of students at the Electrical Engineering and Informatics Faculty, Budapest University of Technology and Economics. The usage of this skeleton is optional. It has been tested with the \emph{TeXLive} \TeX~implementation, and it requires the PDF-\LaTeX~compiler.
\vfill

%----------------------------------------------------------------------------
\chapter*{Bevezet�}\addcontentsline{toc}{chapter}{Bevezet�}
%----------------------------------------------------------------------------

A bevezet� tartalmazza a diplomaterv-ki�r�s elemz�s�t, t�rt�nelmi el�zm�nyeit, a feladat indokolts�g�t (a motiv�ci� le�r�s�t), az eddigi megold�sokat, �s ennek t�kr�ben a hallgat� megold�s�nak �sszefoglal�s�t.

A bevezet� szok�s szerint a diplomaterv fel�p�t�s�vel z�r�dik, azaz annak r�vid le�r�s�val, hogy melyik fejezet mivel foglalkozik.

\chapter*{Els� fejezet}

%----------------------------------------------------------------------------
\chapter*{K�sz�netnyilv�n�t�s}\addcontentsline{toc}{chapter}{K�sz�netnyilv�n�t�s}
%----------------------------------------------------------------------------

Ez nem k�telez�, ak�r t�r�lhet� is. Ha a szerz� sz�ks�g�t �rzi, itt lehet k�sz�netet nyilv�n�tani azoknak, akik hozz�j�rultak munk�jukkal ahhoz, hogy a hallgat� a szakdolgozatban vagy diplomamunk�ban le�rt feladatokat sikeresen elv�gezze. A konzulensnek val� k�sz�netnyilv�n�t�s sem k�telez�, a konzulensnek hivatalosan is dolga, hogy a hallgat�t konzult�lja.


%\listoffigures\addcontentsline{toc}{chapter}{�br�k jegyz�ke}
%\listoftables\addcontentsline{toc}{chapter}{T�bl�zatok jegyz�ke}

\bibliography{mybib}
\addcontentsline{toc}{chapter}{Irodalomjegyz�k}
\bibliographystyle{plain}

%----------------------------------------------------------------------------
\appendix
%----------------------------------------------------------------------------


\label{page:last}
\end{document}
